\subsection{Trigger, Funktionen und Prozeduren}
\subsubsection{Trigger}
Ein \emph{Trigger} fängt ein Ereignis ab (\icode{INSERT}, \icode{UPDATE} oder \icode{DELETE} an einer Tabelle) und führt automatisch einen Anweisungsblock aus („feuern"). Er wird nie vom Benutzer aufgerufen, sondern vom DBMS -- entweder \icode{BEFORE} oder \icode{AFTER} der Änderung. Typische Einsatzzwecke: Werte automatisch setzen, Plausibilität prüfen, Änderungen protokollieren und referenzielle Integrität überwachen.

\subsubsection{Funktionen}
Eine benutzerdefinierte Funktion ist eine \emph{Skalarfunktion}: Sie gibt genau einen Wert zurück. Der Rückgabetyp wird mit \icode{RETURNS} deklariert, der Wert mit \icode{RETURN} zurückgegeben. Nach \icode{DETERMINISTIC} folgt der eigentliche Funktionskörper.

\begin{syntaxbox}[SQL]
CREATE FUNCTION multipliziere(param1 FLOAT(8,2), param2 FLOAT(8,2))
RETURNS FLOAT(8,2)
READS SQL DATA
DETERMINISTIC
RETURN param1 * param2;

SELECT multipliziere(Kosten, 1.19) FROM Leistung;   -- Aufruf
\end{syntaxbox}

\subsubsection{Prozeduren (Stored Procedures)}
Eine \emph{Stored Procedure} ist ein benannter, in der Datenbank gespeicherter Anweisungsblock -- mit oder ohne Parameter, aufgerufen per \icode{CALL}. Nutzen: Businesslogik in die Datenbank auslagern, Wiederverwendbarkeit von mehreren Rechnern aus und Parameterübergabe.
